\documentclass[11pt]{article}

\usepackage[margin=0.82in]{geometry}
\usepackage[T1]{fontenc}
\usepackage{lmodern}
\usepackage{microtype}
\usepackage{amsmath,amssymb,bm,mathtools}
\usepackage{booktabs,tabularx,array}
\usepackage[ruled,vlined,linesnumbered]{algorithm2e}
\usepackage[dvipsnames]{xcolor}
\usepackage{enumitem}
\usepackage{natbib}
\usepackage{hyperref}
\usepackage[nameinlink,noabbrev]{cleveref}

\newcommand{\method}[1]{\textsc{#1}}
\newcommand{\ours}{\textsc{P-Online}}
\newcommand{\pdpo}{\textsc{P-DPO}}
\newcommand{\prlhf}{\textsc{P-RLHF}}
\newcommand{\agenticdpo}{\textsc{Agentic-DPO}}
\newcommand{\shorthand}{\textsc{SHR}}
\newcommand{\profileswap}{\textsc{PSPG}}
\newcommand{\tomorrow}{\textsc{TMNW}}
\newcommand{\forkquery}{\textsc{CFQ}}

\newcommand{\shared}{\bm{s}}
\newcommand{\private}{\bm{p}}
\newcommand{\params}{\bm{\theta}}
\newcommand{\policy}{\pi_{\params}}
\newcommand{\refpolicy}{\pi_{\mathrm{ref}}}
\newcommand{\behpolicy}{\mu}
\newcommand{\loss}{\mathcal{L}}
\newcommand{\buffer}{\mathcal{R}}
\newcommand{\anchors}{\mathcal{A}}
\newcommand{\users}{\mathcal{U}}
\newcommand{\expect}{\mathbb{E}}
\newcommand{\ind}{\mathbb{I}}
\newcommand{\R}{\mathbb{R}}
\newcommand{\norm}[1]{\left\lVert #1 \right\rVert}
\newcommand{\given}{\,\middle|\,}
\newcommand{\softmax}{\operatorname{softmax}}
\newcommand{\clip}{\operatorname{clip}}
\newcommand{\rank}{\operatorname{rank}}
\newcommand{\KL}{D_{\mathrm{KL}}}
\newcommand{\JS}{D_{\mathrm{JS}}}

\newcommand{\Inherited}{\textcolor{MidnightBlue}{\textbf{Inherited}}}
\newcommand{\Proposed}{\textcolor{BrickRed}{\textbf{Proposed}}}
\newcommand{\Evidence}{\textcolor{ForestGreen}{\textbf{Evidence boundary}}}
\newcommand{\code}[1]{\texttt{#1}}

\newcolumntype{Y}{>{\raggedright\arraybackslash}X}
\newcolumntype{L}[1]{>{\raggedright\arraybackslash}p{#1}}

\setlist[itemize]{leftmargin=1.35em,itemsep=2pt,topsep=3pt}
\setlist[enumerate]{leftmargin=1.6em,itemsep=3pt,topsep=3pt}
\setlength{\tabcolsep}{5pt}
\renewcommand{\arraystretch}{1.12}
\SetAlgoNlRelativeSize{-1}
\DontPrintSemicolon


\hypersetup{
  colorlinks=true,
  linkcolor=MidnightBlue,
  citecolor=ForestGreen,
  urlcolor=BrickRed,
  pdftitle={Five Algorithmic Workflows for Offline-to-Online Personalized LLM Training}
}

\title{\bfseries Five Algorithmic Workflows for\\
Offline-to-Online Personalized LLM Training}
\author{Technical Design Note}
\date{\today}

\begin{document}
\maketitle

\begin{abstract}
This note turns five research ideas into executable algorithm specifications.
Each section gives the supporting papers, the inherited method boundary, a
step-by-step workflow, a mathematical objective, pseudocode, invariants, and
rollback conditions.  The common setting starts from an offline personalized
policy and adapts a frozen-backbone language-model agent through shared LoRA
parameters and user-conditioned embeddings.  The five proposals address
ownership-aware stability, stale-trajectory reuse, user-specific policy
advantages, temporal negative weighting, and active preference acquisition.
They are designs for implementation and evaluation, not reports of completed
experiments.
\end{abstract}

\tableofcontents
\newpage

\section{Common Training Contract}
\label{sec:common}

\subsection{Scope and evidence}

All five ideas begin with an offline personalized language-model policy and
continue learning from chronological agent interactions.  The common policy
state is
\begin{equation}
  \params = \bigl(\phi,e_0,\{e_u:u\in\users\}\bigr),
  \label{eq:common-state}
\end{equation}
where the backbone is frozen, $\phi$ is a shared LoRA adapter, $e_0$ is a
generic soft prompt, and $e_u$ is a private embedding for user $u$.
This organization follows personalized RLHF and parameter-efficient
fine-tuning \citep{li2024personalized,hu2022lora}.  Preference losses use DPO
when a pairwise label is available \citep{rafailov2023dpo}; numerical updates
use AdamW \citep{loshchilov2019adamw}.

\Evidence.  The cited papers establish prior algorithms and motivate the
proposed transformations.  They do not establish that any method in this note
works.  Every improvement claim therefore remains a hypothesis until a
checkpoint-changing run and untouched future evaluation have been completed.

\subsection{Shared record and checkpoint schema}

Every trainable record must identify the user and profile version, task and
environment, generating model, behavior checkpoint, rendered state, actions
and tool calls, outcome or preference provenance, timestamp, data split, and
training eligibility.  Token log-probabilities and replayable state
fingerprints are required when an algorithm uses them.  Missing fields cause
quarantine; they are never silently reconstructed.

Each candidate checkpoint belongs to an append-only directed acyclic graph.
Its receipt contains the parent hash, adapter hash, optimizer-state hash,
data-policy version, evaluation receipt, activation status, and rollback
reason.  A rejected candidate cannot generate later data or donate optimizer
moments.

\subsection{Shared lifecycle}

\begin{table}[t]
  \centering
  \caption{Lifecycle used by all five workflows.  The middle step changes by
  idea; the provenance and activation steps do not.}
  \label{tab:shared-lifecycle}
  \begin{tabularx}{\textwidth}{@{}L{0.09\textwidth}Y Y@{}}
    \toprule
    Phase & Required action & Output \\
    \midrule
    Offline & Train the registered \prlhf\ or \agenticdpo\ parent on a
    chronological historical split. & Immutable checkpoint $c_0$, fixed
    reference policy, optimizer receipt. \\
    Collect & Run only the active checkpoint in the declared environment and
    validate every new record. & Versioned interaction ledger. \\
    Adapt & Execute the idea-specific update without reading validation or
    future-test evidence. & Candidate $c_{t+1}$ and complete update receipt. \\
    Verify & Load $c_{t+1}$ in a fresh worker and compare it with $c_t$ on a
    fixed validation panel. & Activation or rollback decision. \\
    Evaluate & Test activated checkpoints on untouched future periods from
    natural initial states. & Separate task-success and personalization
    estimates. \\
    \bottomrule
  \end{tabularx}
\end{table}

\subsection{Upper, middle, and lower algorithm layers}

The \emph{upper layer} names the optimized object: the shared LoRA and
user-conditioned parameters in \cref{eq:common-state}.  The \emph{middle
layer} maps evidence to an update, such as preference optimization, replay,
policy gradient, bilevel weighting, or active acquisition.  The \emph{lower
layer} provides loops, buffers, projections, checkpointing, and stopping
rules.  A method is credited only for the layer it changes.  Standard
minibatching, AdamW, replay storage, and checkpoint hashing remain
infrastructure.

\subsection{Common activation gate}
\label{sec:activation-gate}

A candidate may become active only if all of the following hold:
\begin{enumerate}
  \item its adapter and optimizer hashes load in a clean worker;
  \item the update is finite and nonzero;
  \item task success and safety are non-inferior on the fixed validation panel;
  \item the declared personalization validation statistic improves;
  \item no protected provenance or future-test boundary was crossed.
\end{enumerate}
Failure restores both parent parameters and parent optimizer state.  Runs stop
when their fixed interaction, token, optimizer-step, GPU, or wall-clock budget
is exhausted; after three consecutive windows without an activated candidate;
or at a preregistered safety termination.

\section{Idea 1: Ownership-Cone Personalized DPO}
\label{sec:ownership-cone}

\subsection{Thesis and novelty boundary}

A shared parameter update can change behavior for every user, while an update
to $e_u$ affects only user $u$.  \method{Ownership-Cone P-DPO} matches each
stability constraint to this ownership boundary: other-user replay constrains
the shared block, and same-user replay constrains the active private block.

\Inherited: the \prlhf\ parameterization and specific-plus-generic preference
loss, DPO, LoRA, AdamW, replay, and gradient projection.

\Proposed: project the \emph{actual AdamW displacement} in two separate
ownership-scoped cones.  No global no-forgetting guarantee is claimed.  The
constraint controls only the first-order change of the sampled replay losses.

\begin{table}[t]
  \centering
  \caption{Algorithm papers supporting Ownership-Cone P-DPO.}
  \label{tab:ownership-support}
  \begin{tabularx}{\textwidth}{@{}L{0.24\textwidth}Y Y@{}}
    \toprule
    Supporting paper & Component used here & Missing component supplied by this proposal \\
    \midrule
    P-RLHF \citep{li2024personalized} &
    Shared LoRA, generic embedding, user embeddings, and personalized DPO. &
    It is offline and has no ownership-scoped online stability rule. \\
    DPO, LoRA, AdamW \citep{rafailov2023dpo,hu2022lora,loshchilov2019adamw} &
    Pairwise loss, parameter-efficient update space, and optimizer. &
    They do not decide which users may constrain which parameter block. \\
    GEM and A-GEM \citep{lopezpaz2017gem,chaudhry2019agem} &
    Replay-gradient constraints and projection. &
    Their task constraints do not distinguish shared from user-private
    parameters or the displacement applied by AdamW. \\
    CAGrad \citep{liu2021cagrad} &
    Gradient-conflict control across objectives. &
    It does not encode parameter ownership or chronological personalization. \\
    FedPDPO \citep{zhu2026fedpdpo} &
    Shared and personalized preference optimization. &
    Its federated split does not provide the proposed deployed replay cones. \\
    \bottomrule
  \end{tabularx}
\end{table}

\subsection{Step-by-step workflow}

\begin{enumerate}
  \item Train an immutable offline \prlhf\ parent $c_0$ and freeze
  $\refpolicy$.
  \item Collect preference pairs with the active checkpoint.  Reject pairs
  without a comparator, user identity, behavior hash, and ranking provenance.
  \item Insert eligible pairs into a replay buffer stratified by user, task,
  time, and policy version.
  \item Draw a current minibatch for user $u$, balanced other-user anchors for
  shared parameters, and same-user anchors for $e_u$.
  \item Advance AdamW moments functionally and form the candidate descent
  displacements $d_s$ and $d_p$, including weight decay.
  \item Project $d_s$ through cross-user constraints and $d_p$ through
  same-user constraints.  Keep every $e_v$, $v\ne u$, frozen.
  \item Apply a finite nonzero projected step and write its gradients,
  optimizer state, active constraints, and projection angles to the receipt.
  \item Hash and load the candidate in a clean worker.  Activate it only
  through the gate in \cref{sec:activation-gate}; otherwise restore the parent.
  \item Evaluate activated checkpoints on untouched future periods.  These
  outcomes do not return to the reported training run.
\end{enumerate}

\subsection{Objective and projected update}

For $z=(u,x,y^+,y^-)$ and conditioning embedding $e$, define
\begin{align}
  \Delta q_{\params}(z,e)
  &=
  \log\frac{\pi_{\params}(y^+\mid x,e)}
            {\refpolicy(y^+\mid x,e)}
  -
  \log\frac{\pi_{\params}(y^-\mid x,e)}
            {\refpolicy(y^-\mid x,e)}, \\
  \ell_{\pdpo}(z;\params)
  &=
  -\log\sigma\!\left(\beta\Delta q_{\params}(z,e_u)\right)
  -\lambda_g\log\sigma\!\left(\beta\Delta q_{\params}(z,e_0)\right).
  \label{eq:ownership-pdpo}
\end{align}
Let $\shared=(\phi,e_0)$ and $\private=e_u$.  AdamW proposes descent
displacements $d_s$ and $d_p$ from the current minibatch.  For cross-user
shared anchors $A^s_j$ and same-user private anchors $A^p_k$, set
\begin{equation}
  a^s_j=\nabla_{\shared}\ell_{\pdpo}(A^s_j;\params),
  \qquad
  a^p_k=\nabla_{\private}\ell_{\pdpo}(A^p_k;\params).
\end{equation}
The two projections are
\begin{align}
  d_s^\star
  &=
  \arg\min_d \frac{1}{2}\norm{d-d_s}_2^2
  &&\text{s.t.}\quad
  (a^s_j)^\top d\ge -\epsilon^s_j/\eta
  \quad \forall j, \label{eq:shared-cone}\\
  d_p^\star
  &=
  \arg\min_d \frac{1}{2}\norm{d-d_p}_2^2
  &&\text{s.t.}\quad
  (a^p_k)^\top d\ge -\epsilon^p_k/\eta
  \quad \forall k. \label{eq:private-cone}
\end{align}
Apply $\shared\leftarrow\shared-\eta d_s^\star$ and
$e_u\leftarrow e_u-\eta d_p^\star$.  Since
$\ell(\shared-\eta d)\approx\ell(\shared)-\eta a^\top d$, each constraint
bounds the sampled anchor's first-order loss increase by $\epsilon$.  It says
nothing about higher-order change, unobserved users, or future reward.

\begin{algorithm}[t]
  \caption{Ownership-Cone P-DPO}
  \label{alg:ownership-cone}
  \KwIn{Offline parent $c_0$; replay buffer $\buffer$; validation panel $V$;
  budgets $B$; tolerances $\epsilon^s,\epsilon^p$}
  \KwOut{Append-only sequence of activated checkpoints}
  $c\leftarrow c_0$; freeze backbone and $\refpolicy$\;
  \While{budget and patience remain}{
    collect eligible pairs with behavior checkpoint $c$ and append to $\buffer$\;
    sample active user $u$ and current minibatch $D_u$\;
    sample balanced cross-user anchors $\anchors_s$ and same-user anchors $\anchors_p$\;
    \For{$k=1,\ldots,K$}{
      form actual AdamW displacements $(d_s,d_p)$ on $D_u$\;
      $d_s^\star\leftarrow\operatorname{Project}(d_s,\anchors_s,\epsilon^s)$ using \cref{eq:shared-cone}\;
      $d_p^\star\leftarrow\operatorname{Project}(d_p,\anchors_p,\epsilon^p)$ using \cref{eq:private-cone}\;
      update only $(\phi,e_0,e_u)$; keep $e_v$ fixed for $v\ne u$\;
      log optimizer moments, constraints, norms, and projection angles\;
    }
    save and hash candidate $\hat c$ with parent $c$\;
    \eIf{$\hat c$ loads cleanly and passes $V$}{
      activate $\hat c$; $c\leftarrow\hat c$\;
    }{
      restore parameters and optimizer state of $c$\;
    }
  }
\end{algorithm}

\subsection{Invariants, stopping, and tests}

\begin{itemize}
  \item Other-user anchors never constrain another user's private embedding.
  Same-user anchors never stand in for shared cross-user protection.
  \item A raw-gradient projection is a baseline, not an implementation of the
  proposal; the projected object must be the applied AdamW displacement.
  \item Every candidate has one parent, and failed candidates cannot seed data
  or optimizer state.
  \item Task success and personalization are gated and reported separately.
\end{itemize}

An adaptation window stops at its declared optimizer-step or token budget.
The main attribution tests are unconstrained online \pdpo, one global cone,
shared-only and private-only cones, ownership-swapped cones, and raw-gradient
rather than AdamW-displacement projection.  Equal initialization, pairs,
optimizer settings, trainable parameters, tokens, and validation gates are
required across these controls.

\Evidence.  This section specifies a testable algorithm.  It contains no
measured result and does not imply that the cone is feasible, nonzero, or
beneficial on the target stream.

\section{Idea 2: Support-Handoff Replay}
\label{sec:support-handoff}

\subsection{Thesis and novelty boundary}

A historical action may become unusable for off-policy learning even when the
state it reached remains useful.  \shorthand\ therefore gives each trajectory
a monotone replay status:
\[
  \code{action\_alive}\longrightarrow\code{state\_only}
  \longrightarrow\code{quarantined}.
\]
Supported trajectories contribute importance-weighted action gradients.  Once
support collapses, their actions expire permanently; a verified historical
state may still seed a fresh current-policy continuation.

\Inherited: \prlhf\ initialization, off-policy weights, effective sample size
(ESS), replay buffers, state resets, LoRA, and AdamW.

\Proposed: the irreversible action-to-state handoff and its rule that expired
action tokens can never contribute another gradient.

\begin{table}[t]
  \centering
  \caption{Algorithm papers supporting Support-Handoff Replay.}
  \label{tab:handoff-support}
  \begin{tabularx}{\textwidth}{@{}L{0.24\textwidth}Y Y@{}}
    \toprule
    Supporting paper & Component used here & Missing component supplied by this proposal \\
    \midrule
    P-RLHF \citep{li2024personalized} &
    Offline personalized initialization and trainable parameter blocks. &
    No deployed trajectory-replay loop. \\
    WebRL \citep{qi2025webrl} &
    Online agent RL and replay of useful experience. &
    It does not preserve an expired trajectory solely as an executable state. \\
    VCPO \citep{huang2026stable} &
    ESS as a diagnostic for stale off-policy updates. &
    It stabilizes weighted updates; it does not switch representation from
    action to state. \\
    Balanced Replay \citep{lee2021offlinetoonline} &
    Controlled mixing during offline-to-online transition. &
    It retains or balances samples rather than terminating their action
    eligibility. \\
    Offline/online LLM RL \citep{lanchantin2025bridging} &
    Evidence that synchronization and online data change LLM post-training. &
    It does not define reconstructible state continuation for personalized
    agents. \\
    \bottomrule
  \end{tabularx}
\end{table}

\subsection{Step-by-step workflow}

\begin{enumerate}
  \item Train and hash an offline \prlhf\ checkpoint.
  \item Store each live trajectory twice: an action payload with behavior
  log-probabilities and a state payload with tool-boundary fingerprints.
  \item Group action-alive records by behavior checkpoint and compatible
  reward semantics.
  \item Recompute current-policy log-probabilities and normalized ESS for each
  sampled cohort.
  \item If ESS remains above the fixed threshold, form a clipped,
  self-normalized action-replay loss.
  \item On the first ESS failure, atomically mark the cohort
  \code{state\_only}.  Missing behavior log-probabilities cause the same
  transition.
  \item For a state-only record, replay a sampled prefix in a clean sandbox.
  Continue only if its canonical fingerprint exactly matches the ledger.
  \item Generate fresh suffixes with the active checkpoint and train only on
  suffix tokens using newly verified outcomes.
  \item Mix the available action and state losses, execute AdamW, then apply
  the common candidate-load, activation, and rollback gate.
\end{enumerate}

\subsection{Support test and losses}

For a behavior cohort $C$ of size $n$, define clipped ratios
\begin{equation}
  w_i
  =
  \exp\!\left[
    \clip\!\left(
      \log\policy(\tau_i\mid u_i)-\log\mu_i(\tau_i\mid u_i),
      -\log 10,\log 10
    \right)
  \right]
\end{equation}
and normalized ESS
\begin{equation}
  \eta_C
  =
  \frac{\left(\sum_{i\in C}w_i\right)^2}
       {n\sum_{i\in C}w_i^2}.
  \label{eq:handoff-ess}
\end{equation}
Use action replay only when $n\ge 8$, $\eta_C\ge\kappa$ with
$\kappa=0.5$, and reward semantics remain compatible.  With
$\bar w_i=w_i/\sum_jw_j$ treated as stop-gradient and $A_i$ a return centered
within user-task strata,
\begin{equation}
  \loss_A
  =
  -\sum_{i\in C}\bar w_i A_i
    \sum_{j\in\tau_i}
    \log\policy(a_{ij}\mid h_{ij},e_{u_i}).
  \label{eq:action-replay}
\end{equation}

For a state-only trajectory, sample a verified boundary $s_{ik}$, reconstruct
it, and draw $G=4$ fresh suffixes $\xi_{ig}\sim\policy(\cdot\mid s_{ik},e_u)$.
Using leave-one-out advantages $\widetilde A_{ig}$ from newly verified
rewards,
\begin{equation}
  \loss_S
  =
  -\frac{1}{G}\sum_{g=1}^{G}\widetilde A_{ig}
  \sum_{j\in\xi_{ig}}
  \log\policy(a_{igj}\mid h_{igj},e_u).
  \label{eq:state-replay}
\end{equation}
The update minimizes
\begin{equation}
  \loss
  =
  \frac{1}{2}\loss_A+\frac{1}{2}\loss_S
  +\beta\KL(\policy\Vert\pi_{\mathrm{parent}}),
  \label{eq:handoff-loss}
\end{equation}
renormalizing over the modes available in the minibatch.  The clipped,
self-normalized action estimator is biased by design.

\begin{algorithm}[t]
  \caption{Support-Handoff Replay}
  \label{alg:support-handoff}
  \KwIn{Parent $c_0$; dual-payload ledger $\buffer$; threshold $\kappa$;
  suffix count $G$; validation panel $V$}
  \KwOut{Activated checkpoints and monotone replay statuses}
  $c\leftarrow c_0$\;
  \While{budget and patience remain}{
    collect trajectories with $c$; validate action and state payloads\;
    sample behavior cohorts $\{C\}$ from $\buffer$\;
    $\loss_A\leftarrow 0$; $\loss_S\leftarrow 0$\;
    \ForEach{cohort $C$}{
      compute $\eta_C$ by \cref{eq:handoff-ess}\;
      \eIf{$C$ is \code{action\_alive}, $\eta_C\ge\kappa$, and rewards are compatible}{
        add the loss in \cref{eq:action-replay} to $\loss_A$\;
      }{
        mark $C$ \code{state\_only} permanently\;
        sample a tool boundary and reconstruct it in a clean sandbox\;
        \eIf{the canonical fingerprint matches}{
          generate $G$ suffixes with $c$ and add \cref{eq:state-replay} to $\loss_S$\;
        }{
          mark the record \code{quarantined}\;
        }
      }
    }
    take one clipped AdamW step on the available terms of \cref{eq:handoff-loss}\;
    save and hash candidate $\hat c$\;
    \eIf{$\hat c$ loads and passes $V$}{
      activate $\hat c$; $c\leftarrow\hat c$\;
    }{
      restore $c$ and its optimizer state\;
    }
  }
\end{algorithm}

\subsection{Invariants, stopping, and tests}

\begin{itemize}
  \item Replay status is monotone.  A state-only cohort never returns to action
  replay, even if a later policy raises its apparent ESS.
  \item An expired prefix receives no gradient.  Only fresh suffix tokens can
  train the policy after handoff.
  \item State replay requires a clean-sandbox reconstruction and exact
  canonical fingerprint.  Failure means quarantine, not approximate reuse.
  \item Evaluation always starts from natural task initial states, never from a
  privileged replay boundary.
\end{itemize}

An adaptation window stops after 200 optimizer steps, 500 fresh suffixes, or
its token ceiling.  The necessary controls are fresh-only learning, a fixed
offline/online replay mixture, ESS-based discard, ESS learning-rate scaling,
always-action replay, always-state replay, and a compute-matched sham handoff.
All controls receive the same ledger, reward semantics, suffix budget, and
optimizer implementation.

\Evidence.  The state adapter, reconstruction contract, and learner remain to
be implemented.  Repository readiness or a valid fingerprint alone does not
show that handoff improves either personalization or task success.

\section{Idea 3: Profile-Swap Policy Gradient}
\label{sec:profile-swap}

\subsection{Thesis and novelty boundary}

A trajectory can receive high personalized reward because it is simply good
for everyone.  \profileswap\ holds the realized trajectories fixed and changes
only the profile used for scoring.  Its personalization advantage is the
intended user's within-group rank minus the mean rank under matched alternative
profiles.

\Inherited: the personalized parameterization, grouped policy gradient, frozen
profile scorers, rank normalization, LoRA, and AdamW.

\Proposed: a cross-profile rank contrast on the \emph{same} trajectory group.
It is an evaluative contrast, not a causal estimate of user identity.

\begin{table}[t]
  \centering
  \caption{Algorithm papers supporting Profile-Swap Policy Gradient.}
  \label{tab:profile-support}
  \begin{tabularx}{\textwidth}{@{}L{0.24\textwidth}Y Y@{}}
    \toprule
    Supporting paper & Component used here & Missing component supplied by this proposal \\
    \midrule
    P-RLHF \citep{li2024personalized} &
    Shared LoRA and user-conditioned embeddings learned from preferences. &
    It has no chronological on-policy policy-gradient loop. \\
    DeepSeekMath/GRPO \citep{shao2024deepseekmath} &
    Group-relative advantages and clipped policy optimization. &
    Group normalization does not separate intended-user value from broadly
    useful behavior. \\
    WebRL \citep{qi2025webrl} &
    Online reinforcement learning for multi-turn language-model agents. &
    It has no persistent cross-profile personalization contrast. \\
    Reward factorization \citep{shenfeld2025reward} &
    User-conditioned reward structure and low-dimensional preference
    variation. &
    It does not form an on-policy advantage by rescoring identical
    trajectories under other users. \\
    \bottomrule
  \end{tabularx}
\end{table}

\subsection{Step-by-step workflow}

\begin{enumerate}
  \item Train an offline \prlhf\ checkpoint and freeze the preference scorer
  version used by the next adaptation window.
  \item For user-task context $(u,x)$, generate $G=8$ trajectories with the
  exact active checkpoint.
  \item Before observing scores, sample $K=4$ alternative profiles matched on
  task feasibility, permissions, language, and available evidence.
  \item Score the same eight trajectories under $u$ and every alternative
  profile.  Do not rerun the environment.
  \item Convert each profile's scores into centered within-group ranks.
  \item Subtract the alternative-profile mean rank from the intended-user
  rank.  Separately standardize deterministic task success.
  \item Combine the two advantages with a validation-selected coefficient and
  execute one clipped on-policy AdamW update.
  \item Hash, load, validate, activate, or roll back the candidate.
  \item Evaluate future task success and personalization separately.
\end{enumerate}

\subsection{Profile-swap advantage}

Let $h_\omega(v,x,\tau_i)$ be a frozen score for profile $v$.  For each
$v\in\{u\}\cup C(u,x)$, rank the $G$ trajectories from least to most
preferred, use average ranks for ties, and define
\begin{equation}
  q_{v,i}
  =
  \frac{2\rank_v(i)-G-1}{G-1}.
  \label{eq:profile-rank}
\end{equation}
Thus $q_{v,i}\in[-1,1]$ and $\sum_iq_{v,i}=0$.  The profile-swap advantage is
\begin{equation}
  A^{\mathrm{swap}}_i
  =
  q_{u,i}
  -
  \frac{1}{K}\sum_{v\in C(u,x)}q_{v,i}.
  \label{eq:swap-advantage}
\end{equation}
For deterministic task success $s_i$, define
\begin{equation}
  A^{\mathrm{task}}_i
  =
  \frac{s_i-\bar s}{\operatorname{sd}(s_1,\ldots,s_G)+10^{-6}},
\end{equation}
and set it to zero when the group variance is zero.  The training advantage is
\begin{equation}
  A_i
  =
  \clip(A^{\mathrm{task}}_i,-2,2)
  +
  \lambda\clip(A^{\mathrm{swap}}_i,-2,2),
  \label{eq:profile-total-adv}
\end{equation}
where $\lambda$ is selected on offline validation and then frozen.

If $\mu$ is the behavior policy and
$\rho_{it}=\policy(a_{it}\mid h_{it},u)/\mu(a_{it}\mid h_{it},u)$, minimize
\begin{equation}
  \loss_{\mathrm{PSPG}}
  =
  -\expect_{i,t}\!\left[
    \min\!\left\{
      \rho_{it}A_i,\,
      \clip(\rho_{it},1-\epsilon,1+\epsilon)A_i
    \right\}
  \right]
  +
  \beta\KL(\policy\Vert\pi_{\mathrm{parent}}).
  \label{eq:profile-policy-loss}
\end{equation}
Ranks make the update invariant to a strictly monotone recalibration applied
separately to each profile scorer.

\begin{algorithm}[t]
  \caption{Profile-Swap Policy Gradient}
  \label{alg:profile-swap}
  \KwIn{Parent $c_0$; frozen scorer $h_\omega$; panel sampler; $G,K,\lambda$;
  validation panel $V$}
  \KwOut{Activated personalized policy checkpoints}
  $c\leftarrow c_0$\;
  \While{budget and patience remain}{
    sample eligible user-task context $(u,x)$\;
    generate $\{\tau_i\}_{i=1}^{G}$ with active checkpoint $c$\;
    sample matched alternative panel $C(u,x)$ before scoring\;
    \ForEach{$v\in\{u\}\cup C(u,x)$}{
      score the identical trajectories with $h_\omega(v,x,\cdot)$\;
      compute centered ranks $q_{v,i}$ by \cref{eq:profile-rank}\;
    }
    compute $A_i^{\mathrm{swap}}$, $A_i^{\mathrm{task}}$, and $A_i$\;
    consume the rollout group once in a clipped update on \cref{eq:profile-policy-loss}\;
    update only $(\phi,e_0,e_u)$ and log scorer, panel, policy, and optimizer versions\;
    save and hash candidate $\hat c$\;
    \eIf{$\hat c$ loads and passes $V$}{
      activate $\hat c$; $c\leftarrow\hat c$\;
    }{
      restore $c$ and its optimizer state\;
    }
  }
\end{algorithm}

\subsection{Invariants, stopping, and tests}

\begin{itemize}
  \item Every profile scores the same trajectory tokens; profile swapping
  never changes environment execution.
  \item Alternative profiles are sampled before scores are visible, and their
  identities and propensities are recorded.
  \item The scorer is frozen within an adaptation window and receives no
  policy-training gradient.
  \item Each rollout group is consumed for at most one optimizer epoch.
  \item Task success, safety, and permissions are not counterfactually
  swapped.
\end{itemize}

An adaptation window stops after 200 optimizer steps, 400 rollout groups, or
its token ceiling.  Required controls include raw-reward GRPO, intended-profile
rank only, user-history normalization, global normalization, same-profile
duplicate panels, randomly permuted identities, and a compute-matched extra
scoring baseline.  These tests distinguish identity contrast from rank
normalization and additional scorer calls.

\Evidence.  Cross-profile subtraction may remove generic quality, but it may
also amplify scorer bias or penalize preferences shared by many users.  Only
matched future evaluation can decide whether the resulting policy is more
personalized without losing task competence.

\section{Idea 4: Tomorrow's Mistake}
\label{sec:tomorrows-mistake}

\subsection{Thesis and novelty boundary}

\agenticdpo\ contrasts an accepted action with a plausible same-state mistake.
The hardest current mistake, however, may teach little about the user's next
tasks.  \tomorrow\ weights current mistakes by a one-step estimate of forward
transfer: would correcting this mistake reduce loss on the same user's next
already-observed interaction block?

\Inherited: \agenticdpo's state-level preference loss, SFT anchor, schema
preservation, and negative refresh; \prlhf's user-conditioned parameters;
bilevel differentiation; and AdamW.

\Proposed: chronological support/query blocks and a virtual optimizer step
that assigns weight to each current negative according to next-block transfer.
The one-step hypergradient is local and does not guarantee future reward.

\begin{table}[t]
  \centering
  \caption{Algorithm papers supporting Tomorrow's Mistake.}
  \label{tab:tomorrow-support}
  \begin{tabularx}{\textwidth}{@{}L{0.24\textwidth}Y Y@{}}
    \toprule
    Supporting paper & Component used here & Missing component supplied by this proposal \\
    \midrule
    Agentic-DPO \citep{chen2026agenticdpo} &
    Same-state action preferences, current-policy negatives, SFT anchoring,
    and schema-preserving views. &
    Its hardest-negative rule is offline, non-personalized, and myopic. \\
    P-RLHF \citep{li2024personalized} &
    Shared and user-private parameterization. &
    It has no action-level negative refresh or temporal meta-weighting. \\
    Example reweighting \citep{ren2018learning} &
    A virtual update that scores training examples by a later objective. &
    The query set is generic validation data rather than the same user's next
    chronological interaction block. \\
    Bilevel optimization \citep{franceschi2018bilevel} &
    Differentiation through an inner optimization step. &
    It does not specify agent mistakes, preference losses, or online user
    chronology. \\
    \bottomrule
  \end{tabularx}
\end{table}

\subsection{Step-by-step workflow}

\begin{enumerate}
  \item Run the \agenticdpo\ SFT warm-up, freeze its reference policy, and
  initialize the shared LoRA and user embeddings from offline history.
  \item Partition each user's eligible stream into immutable consecutive
  blocks $B(u,1),B(u,2),\ldots$.
  \item Use $S=B(u,t-1)$ as support and $Q=B(u,t)$ as the meta-query.  Both
  must be observed before the update; $Q$ is not validation or test data.
  \item At every accepted state in $S$ and $Q$, sample four one-step actions
  from the active checkpoint.  Remove duplicates, schema violations, and
  ineligible negatives.
  \item Assign transient logits to support negatives and take one
  differentiable functional AdamW step on their weighted loss.
  \item Evaluate the inherited \agenticdpo\ loss on $Q$ through the virtual
  parameters, and differentiate it with respect to the support logits.
  \item Convert the negative hypergradient into within-state weights.  Delete
  the virtual parameters and moments.
  \item From the original parent state, execute one real weighted
  \agenticdpo\ AdamW step.
  \item Hash, load, validate, activate, or roll back the candidate.  If
  activated, move $Q$ into the support role for the next block.
\end{enumerate}

\subsection{Temporal bilevel objective}

At support state $s_i$, let $a_i^+$ be the accepted action and
$a_{ik}^-$, $k=1,\ldots,K$, be fresh schema-valid negatives.  Let
$\ell_{\mathrm{ADPO}}(i,k;\params)$ denote the inherited anchored,
action-level \agenticdpo\ loss.  Introduce transient logits
$c_i\in\R^K$ and weights
\begin{equation}
  w_{ik}(c_i)=\softmax(c_i/\tau)_k.
\end{equation}
The weighted support loss is
\begin{equation}
  \loss_S(\params,c)
  =
  \frac{1}{|S|}
  \sum_{i\in S}\sum_{k=1}^{K}
  w_{ik}(c_i)\ell_{\mathrm{ADPO}}(i,k;\params).
  \label{eq:tomorrow-support}
\end{equation}

Let the current AdamW state be $o=(m,v,r)$.  A functional step computes
\begin{align}
  g(c) &= \nabla_{\params}\loss_S(\params,c), \\
  m'(c) &= \beta_1m+(1-\beta_1)g(c), \\
  v'(c) &= \beta_2v+(1-\beta_2)g(c)^2, \\
  \widetilde\params(c)
  &=
  \params-\alpha\left[
    \frac{\widehat m'(c)}{\sqrt{\widehat v'(c)}+\varepsilon}
    +\lambda_{\mathrm{wd}}\params
  \right].
  \label{eq:tomorrow-virtual}
\end{align}
Using uniform query-negative weights, define
\begin{equation}
  J(c)
  =
  \frac{1}{|Q|K}
  \sum_{q\in Q}\sum_{k=1}^{K}
  \ell_{\mathrm{ADPO}}\!\left(q,k;\widetilde\params(c)\right).
  \label{eq:tomorrow-query}
\end{equation}
Starting from $c=0$, take one meta step,
\begin{equation}
  c^\star=-\eta_{\mathrm{meta}}\nabla_cJ(c)\big|_{c=0},
  \qquad
  w^\star_{ik}=\softmax(c_i^\star/\tau)_k.
  \label{eq:tomorrow-weights}
\end{equation}
Then discard $\widetilde\params,m',v'$, stop gradients through $w^\star$, and
take one real AdamW step on $\loss_S(\params,w^\star)$ from the unmodified
parent state.

\begin{algorithm}[t]
  \caption{Tomorrow's Mistake temporal meta-weighting}
  \label{alg:tomorrows-mistake}
  \KwIn{Parent $c_0$; ordered user blocks; negatives per state $K$;
  meta rate $\eta_{\mathrm{meta}}$; validation panel $V$}
  \KwOut{Activated checkpoints and negative-weight receipts}
  $c\leftarrow c_0$\;
  \While{budget and patience remain}{
    choose user $u$ with adjacent eligible blocks $S=B(u,t-1)$ and $Q=B(u,t)$\;
    \If{$S$ or $Q$ lacks the minimum eligible states}{buffer evidence and continue}
    generate $K$ valid same-state negatives in $S$ and $Q$ with checkpoint $c$\;
    initialize support logits $c_i\leftarrow 0$\;
    clone adapter and AdamW state functionally\;
    compute virtual parameters $\widetilde\params(c_i)$ by \cref{eq:tomorrow-virtual}\;
    evaluate $J(c)$ on $Q$ and compute $c^\star$ by \cref{eq:tomorrow-weights}\;
    delete every virtual parameter and moment; restore the bitwise parent state\;
    stop gradients through $w^\star$ and take one real weighted update on $S$\;
    log block, negative-generator, hypergradient, optimizer, and policy hashes\;
    save and hash candidate $\hat c$\;
    \eIf{$\hat c$ loads and passes $V$}{
      activate $\hat c$; $c\leftarrow\hat c$; advance the block pair\;
    }{
      restore $c$ and its optimizer state\;
    }
  }
\end{algorithm}

\subsection{Invariants, stopping, and tests}

\begin{itemize}
  \item Support and meta-query blocks belong to the same authenticated user and
  adjacent ordered periods.  Future blocks cannot be backfilled.
  \item Virtual checkpoints are ephemeral, bitwise isolated, and never
  activatable.
  \item The query block is training evidence already observed before the
  update.  It is disjoint from checkpoint-selection validation and untouched
  future evaluation.
  \item Every negative records its generating checkpoint and is consumed in
  one real update window.
  \item A rejected real candidate supplies neither future rollouts nor AdamW
  moments.
\end{itemize}

Skip a user update when either block has fewer than 16 eligible action states
or no valid negative.  Stop a window after 150 real optimizer steps, 300
negative-generation batches, or its token ceiling.  Controls must include
unmodified \agenticdpo, personalized hardest-negative \agenticdpo, uniform
negative weights, generic validation reweighting, temporally shuffled
same-user queries, other-user queries, and a sham method that pays for the
virtual pass but applies the parent rule.

\Evidence.  A low next-block loss after one virtual step is only a local
compatibility signal.  The experiment must test whether that signal predicts
future personalization and task success, and whether its extra differentiation
cost is justified.

\section{Idea 5: Checkpoint-Fork Querying}
\label{sec:checkpoint-fork}

\subsection{Thesis and novelty boundary}

An uncertain label is not necessarily useful: either answer may induce nearly
the same policy.  \forkquery\ asks for feedback only when the two possible
labels produce virtual one-step checkpoints with different behavior on
user-relevant probe states.

\Inherited: \prlhf\ and \pdpo, active preference collection, AdamW, LoRA,
Jensen--Shannon divergence, and information-gain acquisition.

\Proposed: score a query in \emph{behavior space} after both exact
label-conditioned optimizer branches, then commit only the observed branch.
The finite probe score is not a guarantee of downstream utility.

\begin{table}[t]
  \centering
  \caption{Algorithm papers supporting Checkpoint-Fork Querying.}
  \label{tab:fork-support}
  \begin{tabularx}{\textwidth}{@{}L{0.24\textwidth}Y Y@{}}
    \toprule
    Supporting paper & Component used here & Missing component supplied by this proposal \\
    \midrule
    P-RLHF \citep{li2024personalized} &
    Personalized DPO loss and shared/private parameterization. &
    It trains a fixed preference set and never allocates live queries. \\
    SELM \citep{zhang2024selm} &
    Active exploration for online language-model alignment. &
    It does not value both labels by the behavior of their resulting
    checkpoints. \\
    Hybrid Preference Optimization \citep{bose2024hybrid} &
    Joint use of offline preference data and online queries. &
    It does not define a decision-focused per-query acquisition score. \\
    SPRInG \citep{kim2026spring} &
    Selective updates under continual personalization. &
    Likelihood novelty does not test whether two labels change future actions. \\
    Information-gain acquisition \citep{mackay1992information} &
    Select data by expected reduction of decision uncertainty. &
    Classical information gain does not include a functional LLM optimizer
    branch or agent action probes. \\
    \bottomrule
  \end{tabularx}
\end{table}

\subsection{Step-by-step workflow}

\begin{enumerate}
  \item Train and load an immutable offline \prlhf\ parent.
  \item In each user window, generate a fixed pool of same-task live trajectory
  pairs with the exact active checkpoint.
  \item Keep only pairs matched on deterministic task success and safety, so
  the requested label targets personalization rather than obvious correctness.
  \item Before scoring the pool, freeze one anchor minibatch, user-relevant
  probe states, and schema-valid action slates.
  \item For each candidate and each possible label, clone the parent adapter
  and AdamW state functionally and execute the exact one-step \pdpo\ update.
  \item Evaluate both virtual checkpoints on every fixed action slate.
  \item Compute the label-prior-weighted Jensen--Shannon divergence between
  their action distributions and choose the largest fork.
  \item Abstain if the maximum score is below the calibrated threshold.
  Otherwise request one label and execute only its branch as a real update from
  the unchanged parent.
  \item Hash, load, validate, activate, or roll back the real candidate.  All
  virtual branches and unqueried labels expire.
\end{enumerate}

\subsection{Virtual branches and fork score}

For candidate pair $q=(u,x,y^A,y^B)$, define
\begin{equation}
  d_{\params,e}(q)
  =
  \log\frac{\pi_{\params}(y^A\mid x,e)}
            {\refpolicy(y^A\mid x,e)}
  -
  \log\frac{\pi_{\params}(y^B\mid x,e)}
            {\refpolicy(y^B\mid x,e)}.
\end{equation}
For hypothetical label $z\in\{+1,-1\}$, where $+1$ prefers $y^A$, use
\begin{equation}
  \ell_z(q;\params)
  =
  -\log\sigma\!\left(\beta z d_{\params,e_u}(q)\right)
  -
  \lambda_g\log\sigma\!\left(\beta z d_{\params,e_0}(q)\right).
  \label{eq:fork-pdpo}
\end{equation}
Freeze a seven-pair eligible anchor minibatch $H_t$ for the whole acquisition
window and set
\begin{equation}
  \loss_{q,z}
  =
  \frac{1}{8}\left[
    \ell_z(q;\params)+\sum_{h\in H_t}\ell(h;\params)
  \right].
\end{equation}
From the same parent parameters and AdamW state $(m,v,r)$, each hypothetical
branch computes
\begin{align}
  g_{q,z} &= \nabla_{\params}\loss_{q,z}, \\
  m'_{q,z} &= \beta_1m+(1-\beta_1)g_{q,z}, \\
  v'_{q,z} &= \beta_2v+(1-\beta_2)g_{q,z}^2, \\
  \widetilde\params_{q,z}
  &=
  \params-\alpha\left[
    \frac{\widehat m'_{q,z}}{\sqrt{\widehat v'_{q,z}}+\varepsilon}
    +\lambda_{\mathrm{wd}}\params
  \right].
  \label{eq:fork-branch}
\end{align}

For each of $R=8$ training-side probe states $s_r$, fix a slate
$A_r=\{a_{r1},\ldots,a_{rK}\}$ of $K=4$ valid actions.  The branch
distribution is
\begin{equation}
  P_{q,z,r}(j)
  =
  \softmax_j\!\left(
    |a_{rj}|^{-1}
    \log\pi_{\widetilde\params_{q,z}}(a_{rj}\mid s_r,e_u)
  \right).
\end{equation}
Calibrate the label prior $p_q=\sigma(d_{\params,e_u}(q)/T_c)$ only on
eligible historical labels, and let
$M_{q,r}=p_qP_{q,+,r}+(1-p_q)P_{q,-,r}$.  The acquisition score is
\begin{equation}
  F(q)
  =
  \frac{1}{R}\sum_{r=1}^{R}
  \left[
    p_q\KL(P_{q,+,r}\Vert M_{q,r})
    +(1-p_q)\KL(P_{q,-,r}\Vert M_{q,r})
  \right].
  \label{eq:fork-score}
\end{equation}
This is the mutual information between the binary label and the finite action
slate under the approximation above.

\begin{algorithm}[t]
  \caption{Checkpoint-Fork Querying}
  \label{alg:checkpoint-fork}
  \KwIn{Parent $c_0$; candidate-pool size $M$; anchors $H_t$; probes and
  slates; threshold $\delta$; validation panel $V$}
  \KwOut{Queried labels, abstentions, and activated checkpoints}
  $c\leftarrow c_0$\;
  \While{label budget, compute budget, and patience remain}{
    collect $M$ same-task live pairs with active checkpoint $c$\;
    filter to matched task-success and safety status\;
    freeze $H_t$, probes, slates, and calibrated label priors\;
    \ForEach{candidate $q$}{
      \ForEach{$z\in\{-1,+1\}$}{
        clone the same parent adapter and AdamW state functionally\;
        compute virtual branch $\widetilde\params_{q,z}$ by \cref{eq:fork-branch}\;
        evaluate fixed action slates and discard the virtual state\;
      }
      compute $F(q)$ by \cref{eq:fork-score}\;
    }
    $q^\star\leftarrow\arg\max_qF(q)$\;
    \eIf{$F(q^\star)<\delta$}{
      record an abstention; expire the pool\;
    }{
      request label $z^\star$ only for $q^\star$\;
      from the unchanged parent, take one real AdamW step on $\loss_{q^\star,z^\star}$\;
      save and hash candidate $\hat c$\;
      \eIf{$\hat c$ loads and passes $V$}{
        activate $\hat c$; $c\leftarrow\hat c$\;
      }{
        restore $c$ and its optimizer state\;
      }
    }
  }
\end{algorithm}

\subsection{Invariants, stopping, and tests}

\begin{itemize}
  \item Both virtual labels start from the same parent parameters, moments,
  anchors, probes, learning rate, and weight decay.
  \item Virtual branches leave persistent parameters and moments bitwise
  unchanged and can never be activated.
  \item Labels remain hidden until a candidate has been selected and its query
  receipt written.
  \item Probe states come from eligible training-side history, never validation
  or untouched future tasks.
  \item Every acquisition baseline sees the same pool, hidden labels, feedback
  budget, and real-update implementation.
\end{itemize}

A user window stops after one query and real update, an abstention, or its
rollout/token ceiling.  Required acquisition baselines are random selection,
label entropy, SPRInG-style novelty, expected gradient norm, expected model
change, and a compute-matched sham that executes both branches and probes but
selects randomly.  Parameter-space fork distance and first-order logit
approximations test whether full behavior-space simulation is necessary.

\Evidence.  The method may spend substantial compute to avoid a single weak
label, and its score depends on probe coverage and prior calibration.  The
primary empirical question is personalization per acquired label and per
activated update, with task success, safety, and total GPU cost reported
separately.

\section{The Two Closest Per-User Training Methods}
\label{sec:closest-per-user}

\subsection{What counts as a close match}

Our unit of personalization is one user.  A deployment run owns that user's
adapter, optimizer state, memory, and checkpoint history.  The initial policy
is trained from the user's thoughts, conversations, corrections, and task
records.  Later interactions update the same policy.  A shared backbone may
be trained separately, but the target user's raw history does not train
another user's deployed branch.

Under this definition, the two closest papers are OPPU
\citep{tan2024oppu} and SPRInG \citep{kim2026spring}.  OPPU is the closest
offline initializer: it trains one adapter from one user's history.  SPRInG
is the closest continual method: it carries that user's adapter forward and
updates it from chronological interaction batches.  Together they bracket
our offline-to-online workflow.

\begin{table}[t]
  \centering
  \caption{Why OPPU and SPRInG are the closest matches.}
  \label{tab:closest-match}
  \begin{tabularx}{\textwidth}{@{}L{0.18\textwidth}YYY@{}}
    \toprule
    Property & OPPU & SPRInG & Our setting \\
    \midrule
    Policy owner &
    One PEFT module per user &
    One evolving LoRA adapter per user &
    One user-owned policy branch \\
    Historical input &
    Paired behaviors or plain text &
    Ordered query--response records &
    Thoughts, conversations, task trajectories, and outcomes \\
    Offline update &
    Cross-entropy or next-token training &
    Pretrained period-0 adapter &
    User-history initialization \\
    Later update &
    None &
    Selective SFT on each new period &
    Preference or outcome-driven online update \\
    Persistent output &
    Personal adapter &
    Adapter lineage and bounded memory &
    Adapter, optimizer state, memory, and audit receipt \\
    \bottomrule
  \end{tabularx}
\end{table}

We exclude methods that train one policy across many users and personalize
only through a profile or prompt.  They do not produce the user-owned policy
studied here.

\subsection{OPPU: one history, one adapter}

\paragraph{Model and data.}
OPPU freezes a task-adapted base model with parameters $\theta$ and assigns
user $u$ a private PEFT module $\Delta\theta_u$.  The deployed model is
\begin{equation}
  M_u = M_{\theta\oplus\Delta\theta_u}.
  \label{eq:oppu-model}
\end{equation}
The user's history $H_u^{<t}$ contains only records observed before the
current query.  A record is either a supervised pair $(x_i,y_i)$ or a plain
text sequence $x_i$.  The first form covers prior requests and accepted
responses.  The second can hold notes, conversation text, or other
user-authored material.

\paragraph{Workflow.}
\begin{enumerate}
  \item Adapt a shared base model to the task using users outside the target
  cohort, then merge the task LoRA into the base.
  \item Freeze the shared base.
  \item Initialize one LoRA module $\Delta\theta_u$ for the target user.
  \item Train it only on $H_u$.  When retrieval is enabled, each example may
  retrieve only records that preceded that example.
  \item Save $\Delta\theta_u$.  At inference, load the frozen base and this
  user's adapter, then answer the new query.
\end{enumerate}

\paragraph{Optimization.}
For paired history, OPPU minimizes conditional cross-entropy:
\begin{equation}
  \loss_{\mathrm{pair}}(\Delta\theta_u)
  =
  -\sum_{(x_i,y_i)\in H_u}
   \sum_{j=1}^{|y_i|}
   \log p_{\theta,\Delta\theta_u}
   \left(y_{i,j}\mid \varphi(x_i),y_{i,<j}\right).
  \label{eq:oppu-paired}
\end{equation}
When history contains plain text only, the target is the next token:
\begin{equation}
  \loss_{\mathrm{text}}(\Delta\theta_u)
  =
  -\sum_{x_i\in H_u}
   \sum_{j=1}^{|x_i|}
   \log p_{\theta,\Delta\theta_u}
   \left(x_{i,j}\mid x_{i,<j}\right).
  \label{eq:oppu-text}
\end{equation}
The paper uses rank-8 LoRA on the query and value projections, learning rate
$10^{-5}$, and two or three fixed epochs.  It does not report an online
update, replay rule, rollback rule, or optimizer-state lifecycle.

\begin{algorithm}[t]
  \caption{OPPU for one user}
  \label{alg:oppu-one-user}
  \KwIn{Frozen task model $M_\theta$; chronological user history $H_u$;
  optional retriever $R$}
  \KwOut{Personal adapter $\Delta\theta_u$}
  initialize $\Delta\theta_u$\;
  \ForEach{record $h_i\in H_u$ in chronological order}{
    construct its target from $(x_i,y_i)$ or right-shifted text $x_i$\;
    if retrieval is enabled, retrieve only from $H_u^{<i}$\;
    accumulate the loss in \cref{eq:oppu-paired} or \cref{eq:oppu-text}\;
  }
  update only $\Delta\theta_u$ for a fixed number of epochs\;
  save $\Delta\theta_u$ with its user and history-version identifiers\;
\end{algorithm}

\paragraph{Inputs and outputs.}
Training consumes a frozen base model, one user identifier, and that user's
historical records.  It outputs one loadable adapter.  Inference consumes the
adapter and a new query; optional retrieval or a generated profile may also
enter the prompt.  The paper evaluates Llama-2-7B on seven LaMP tasks using
the 100 users with the longest histories.

\paragraph{What transfers.}
OPPU gives us the correct ownership boundary and a practical offline
initializer.  It does not learn from task outcomes, pairwise preferences, or
live feedback.  It also treats each record as text rather than as a replayable
agent trajectory.  Our offline stage can use its per-user adapter design, but
our record schema must retain actions, tools, outcomes, timestamps, and the
generating checkpoint.

\subsection{SPRInG: one adapter through time}

\paragraph{Model and data.}
SPRInG starts from the same ownership model.  User $u$ has a frozen backbone
$M_\theta$, an adapter $\phi_u^{(t)}$, and a bounded memory
$B_u^{(t)}$.  New interactions arrive in periods:
\begin{equation}
  D_u^{(t)}=\{(q_i,a_i)\}_{i=1}^{n_t}.
  \label{eq:spring-stream}
\end{equation}
Period 0 supplies the initial adapter.  Each later period updates that same
adapter.  The paper uses text-generation records, not tool trajectories or
explicit rewards.

\paragraph{Workflow.}
\begin{enumerate}
  \item Score every new $(q,a)$ under the frozen base and the current
  user-adapted model.
  \item Keep the top 30\% by drift score.
  \item Copy $\phi_u^{(t-1)}$ and fine-tune the copy on the selected records.
  \item Merge selected records with the old memory, score them again under
  the updated adapter, and retain the hardest 50.
  \item For a new query, retrieve two memory records.  Use them only when
  their similarity passes a fixed gate.
  \item Decode by interpolating adapter-only and retrieval-conditioned token
  probabilities.
\end{enumerate}

\paragraph{Selection and update.}
Let
\begin{align}
  p_{\mathrm{base}}(q,a)
  &=P(a\mid q;M_\theta),\\
  p_{\mathrm{adap}}^{(t-1)}(q,a)
  &=P(a\mid q;M_{\theta,\phi_u^{(t-1)}}).
\end{align}
SPRInG scores a record by
\begin{equation}
  S_t(q,a)
  =
  -\log\frac{p_{\mathrm{adap}}^{(t-1)}(q,a)}
                  {p_{\mathrm{base}}(q,a)}
  +\alpha\log p_{\mathrm{base}}(q,a).
  \label{eq:spring-score}
\end{equation}
The likelihood ratio favors behavior not yet captured by the adapter.  The
base likelihood suppresses malformed or low-quality text.  If $\tau_p$ is
the period's percentile threshold,
\begin{equation}
  D_{\mathrm{drift}}^{(t)}
  =
  \{(q,a)\in D_u^{(t)}:S_t(q,a)\ge\tau_p\}.
  \label{eq:spring-drift}
\end{equation}
The adapter then takes supervised next-token updates on
$D_{\mathrm{drift}}^{(t)}$ only:
\begin{equation}
  \phi_u^{(t)}
  =
  \operatorname{SFT}
  \left(\phi_u^{(t-1)},D_{\mathrm{drift}}^{(t)}\right).
  \label{eq:spring-update}
\end{equation}

\paragraph{Residual memory and inference.}
After the update, SPRInG rescans
\begin{equation}
  C_u^{(t)}
  =
  D_{\mathrm{drift}}^{(t)}\cup B_u^{(t-1)}
\end{equation}
and recomputes \cref{eq:spring-score} with $\phi_u^{(t)}$.  It keeps the
$N_{\max}$ highest-scoring records:
\begin{equation}
  B_u^{(t)}
  =
  \operatorname{Top}_{N_{\max}}
  \{(q,a)\in C_u^{(t)};S'_t(q,a)\}.
  \label{eq:spring-buffer}
\end{equation}
For a test query, the same adapted model produces an adapter-only
distribution $p_{\mathrm{par}}$ and, when memory passes the relevance gate, a
retrieval-conditioned distribution $p_{\mathrm{ret}}$.  The final
distribution is
\begin{equation}
  p_{\mathrm{fin}}
  =
  \lambda p_{\mathrm{ret}}+(1-\lambda)p_{\mathrm{par}}.
  \label{eq:spring-interpolation}
\end{equation}
Memory affects inference only.  It is not replayed into later optimizer
steps.

\begin{algorithm}[t]
  \caption{SPRInG for one user}
  \label{alg:spring-one-user}
  \KwIn{Frozen $M_\theta$; initial $\phi_u^{(0)}$; stream
  $\{D_u^{(t)}\}_{t\ge1}$; percentile $p$; budget $N_{\max}$}
  \KwOut{Adapter lineage $\{\phi_u^{(t)}\}$ and memories
  $\{B_u^{(t)}\}$}
  $B_u^{(0)}\leftarrow\varnothing$\;
  \For{$t=1,2,\ldots$}{
    score $D_u^{(t)}$ with \cref{eq:spring-score}\;
    retain $D_{\mathrm{drift}}^{(t)}$ with \cref{eq:spring-drift}\;
    $\phi_u^{(t)}\leftarrow
      \operatorname{SFT}(\phi_u^{(t-1)},D_{\mathrm{drift}}^{(t)})$\;
    $C_u^{(t)}\leftarrow D_{\mathrm{drift}}^{(t)}\cup B_u^{(t-1)}$\;
    rescore $C_u^{(t)}$ under $\phi_u^{(t)}$\;
    retain $B_u^{(t)}$ with \cref{eq:spring-buffer}\;
  }
\end{algorithm}

\paragraph{Inputs and outputs.}
Training consumes one user's initial adapter and a chronological stream of
query--response batches.  Each period outputs a new adapter and memory
snapshot.  Inference consumes the current adapter, current memory, and a new
query, and outputs a response.  The main experiment uses 100 LongLaMP users,
five chronological periods, Gemma-3-4B-IT, rank-8 LoRA, AdamW, two epochs,
and a 90/10 update/evaluation split within each period.

\paragraph{What transfers.}
SPRInG is the closest end-to-end prior because it preserves user ownership
through time.  Its main gap is the learning signal: difficult-to-predict text
is treated as preference drift, and updates still maximize likelihood.  Our
agent receives richer evidence.  A completed task can include tool calls,
environment outcomes, user corrections, and explicit preferences.  Those
signals should determine eligibility and reward; likelihood novelty alone
cannot tell whether an action helped the user.

\subsection{Concrete design for our policy}

The two papers suggest a clean baseline for our candidate online updates:
\begin{enumerate}
  \item Serialize user $u$'s thoughts, conversations, and successful task
  records into a chronological offline history.
  \item Train a user-owned adapter with the OPPU objective.  Save the adapter,
  optimizer state, data boundary, and base-model hash as checkpoint $c_u^{(0)}$.
  \item Collect later tasks only with the active user checkpoint.  Preserve
  states, actions, tools, outcomes, corrections, and behavior log-probabilities.
  \item Use a SPRInG-style selection stage, but score evidence by task
  validity, preference information, and provenance rather than likelihood
  novelty alone.
  \item Apply the chosen update from this note.  Activate it only after the
  shared validation and rollback gate.
\end{enumerate}

This baseline separates inheritance from novelty.  OPPU supplies the offline
per-user policy.  SPRInG supplies the chronological adapter lineage.  Our
contribution begins where supervised text prediction ends: agent evidence,
preference- or outcome-driven updates, auditable optimizer state, and safe
activation.


\clearpage
\bibliographystyle{plainnat}
\bibliography{references}

\end{document}
